% ===========================================================================
% Flux de livraison et de validation.
%
% Relevé dans .github/workflows/ci.yml le 2026-09-22.
%
% Le diagramme doit être honnête sur son point d'arrêt : la chaîne VALIDE mais
% ne PUBLIE aucun artefact et ne DÉPLOIE nulle part. Tout ce qui suit la porte
% de qualité est, à ce jour, manuel et non capturé. Représenter une flèche vers
% un environnement de production serait une invention.
% ===========================================================================
\begin{tikzpicture}[node distance=7mm]

  % ═══ Déclenchement ══════════════════════════════════════════════════════
  \node[pqbox, text width=26mm] (push)
    {Poussée ou\\demande de fusion\\{\scriptsize \texttt{main}, \texttt{master}}};

  % ═══ Chaîne d'intégration — deux travaux parallèles ═════════════════════
  % ─ Travail backend ─
  \node[pqbox, right=13mm of push, yshift=17mm, text width=30mm] (bdeps)
    {Dépendances Python\\{\scriptsize 3.14}};
  \node[pqbox, right=of bdeps, text width=32mm] (bmig)
    {Contrôle de dérive\\du schéma\\{\scriptsize migrations vs modèles}};
  \node[pqbox, right=of bmig, text width=30mm] (btest)
    {\texttt{pytest}\\{\scriptsize PostgreSQL + pgvector,\\ Redis}};

  \begin{scope}[on background layer]
    \node[pqgroup, fit=(bdeps)(bmig)(btest)] (bjob) {};
  \end{scope}
  \node[pqcaption, anchor=south] at (bjob.north) {Travail « Backend — tests »};

  % ─ Travail frontend ─
  \node[pqbox, right=13mm of push, yshift=-20mm, text width=30mm] (fdeps)
    {Dépendances Node};
  \node[pqbox, right=of fdeps, text width=32mm] (fbuild)
    {Build de production\\{\scriptsize typage strict, budgets}};
  \node[pqbox, right=of fbuild, text width=30mm] (ftest)
    {Tests unitaires\\{\scriptsize Chrome sans interface}};

  \begin{scope}[on background layer]
    \node[pqgroup, fit=(fdeps)(fbuild)(ftest)] (fjob) {};
  \end{scope}
  \node[pqcaption, anchor=north] at (fjob.south) {Travail « Frontend — build »};

  % ═══ Porte de qualité ═══════════════════════════════════════════════════
  \node[pqcore, right=13mm of bmig, xshift=34mm, yshift=-17mm,
        text width=26mm] (gate)
    {Porte de\\qualité};

  % ═══ Point d'arrêt ══════════════════════════════════════════════════════
  \node[pqext, right=13mm of gate, text width=32mm] (images)
    {Construction des images\\{\scriptsize \textbf{manuelle}, non capturée}};
  \node[pqext, below=of images, text width=32mm] (deploy)
    {Déploiement\\{\scriptsize \textbf{aucun environnement cible}}};
  \node[pqext, above=of images, text width=32mm] (rollback)
    {Retour arrière\\{\scriptsize \textbf{aucune procédure}}};

  \begin{scope}[on background layer]
    \node[pqgroup, fit=(images)(deploy)(rollback), draw=piqbitWarn] (out) {};
  \end{scope}
  \node[pqcaptionwarn, anchor=south] at (out.north) {Hors chaîne automatisée};

  % ═══ Liens ══════════════════════════════════════════════════════════════
  \draw[pqflow] (push.east) -- (bdeps.west);
  \draw[pqflow] (push.east) -- (fdeps.west);
  \draw[pqflow] (bdeps.east) -- (bmig.west);
  \draw[pqflow] (bmig.east)  -- (btest.west);
  \draw[pqflow] (fdeps.east) -- (fbuild.west);
  \draw[pqflow] (fbuild.east) -- (ftest.west);
  \draw[pqflow] (btest.east) -- (gate.north west);
  \draw[pqflow] (ftest.east) -- (gate.south west);
  \draw[pqopt]  (gate.east)  -- (out.west);

  % ═══ Résultat daté ══════════════════════════════════════════════════════
  % Ancré sous la rangée la PLUS BASSE (travail frontend) et non sous la porte
  % de qualité, qui est à mi-hauteur : la note y recouvrait les étapes du
  % travail frontend et son propre intitulé de groupe.
  \node[pqwarn, below=11mm of fjob.south west, anchor=north west, text width=140mm]
    {\textbf{État capturé le 22 septembre 2026.}
     Build de production : \textsc{réussite} (code 0, 36,0\,s) — mais
     \emph{quatre avertissements de budget}, le bundle initial dépassant la
     cible. Suite backend : \textsc{échec} (594 réussis, 3 échoués, 6 ignorés,
     code 1). La porte de qualité était donc \textbf{fermée} à cette date.
     Une nouvelle campagne datée est requise après correction.};

\end{tikzpicture}
